\projectprogress
\label{sec:progress}

\subsection{Pipeline Architecture}
The project is implemented as a set of Python modules chained into five
sequential stages. Each stage is also exposed as a standalone entry point, so
any part of the workflow can be re-run on its own. All five stages now run end
to end; the completion status is summarized in
Table~\ref{tab:pipeline_status}.

\begin{table}[htbp]
	\centering
	\caption{Pipeline stages and implementation status. Stage~5 is split into
	the auxiliary composition surrogate (5a, Section~\ref{sec:comp_nn}) and
	the parameter-level optimizer (5b, Section~\ref{sec:nnip}).}
	\label{tab:pipeline_status}
	\renewcommand{\arraystretch}{1.2}
	\begin{tabular}{|c|p{6.5cm}|c|c|}
		\hline
		\textbf{Stage} & \textbf{Description} & \textbf{Framework} & \textbf{Status} \\ \hline
		1  & Random alloy configuration generation       & Python                  & Complete \\ \hline
		2  & MD elastic-property evaluation               & LAMMPS Python API       & Complete \\ \hline
		3  & DFT reference calculations                    & Quantum Espresso + ASE  & Complete \\ \hline
		4  & MEAM initialization and library merging       & Python                  & Complete \\ \hline
		5a & Composition surrogate (sanity check)          & TensorFlow              & Complete \\ \hline
		5b & NN parameter optimization (closes the loop)   & TensorFlow + LAMMPS     & Complete \\ \hline
	\end{tabular}
\end{table}

Figure~\ref{fig:pipeline} shows the data flow. Stage~1 emits JSON
descriptions of alloy configurations. Stage~2 simulates them in LAMMPS and
writes the measured elastic constants to a shared results database. Stage~3
produces DFT reference values for every selected element and binary pair.
Stage~4 reads both databases and emits a merged, DFT-initialized MEAM
library. Stage~5b then consumes that library together with the
LAMMPS-evaluated $(C_{11},C_{12})$ targets and produces a refitted parameter
set. The composition-only surrogate (dashed in Figure~\ref{fig:pipeline})
reads the Stage~2 database as a consistency check; it is described in
Section~\ref{sec:comp_nn}.

\begin{figure}[htbp]
	\centering
	\begin{tikzpicture}[
		stage/.style={rectangle, rounded corners, draw=black, fill=gray!10,
		              text width=2.6cm, minimum height=1.0cm, align=center, font=\small\bfseries},
		aux/.style={rectangle, rounded corners, draw=black, dashed, fill=blue!5,
		              text width=2.6cm, minimum height=1.0cm, align=center, font=\small\bfseries},
		arrow/.style={-Stealth, thick},
		auxarrow/.style={-Stealth, thick, dashed},
		label/.style={font=\scriptsize, align=center},
		node distance=0.4cm and 0.4cm,
	]
		\node[stage] (s1) {Stage 1\\Configurations};
		\node[stage, right=of s1] (s2) {Stage 2\\MD: $C_{11},C_{12}$};
		\node[stage, right=of s2] (s3) {Stage 3\\DFT};
		\node[stage, right=of s3] (s4) {Stage 4\\MEAM init};
		\node[stage, right=of s4] (s5) {Stage 5b\\NN optim};

		\draw[arrow] (s1) -- (s2);
		\draw[arrow] (s2) -- (s3);
		\draw[arrow] (s3) -- (s4);
		\draw[arrow] (s4) -- (s5);

		\node[label, below=0.1cm of s1] {composition\\records};
		\node[label, below=0.1cm of s2] {elastic\\database};
		\node[label, below=0.1cm of s3] {DFT\\reference};
		\node[label, below=0.1cm of s4] {merged MEAM\\library};
		\node[label, below=0.1cm of s5] {refitted MEAM\\parameters};

		% Auxiliary composition surrogate: consumes the Stage 2 database
		% as a consistency check, separate from the potential-fitting flow.
		\node[aux, above=0.9cm of s3] (ml) {Stage 5a\\composition NN};
		\draw[auxarrow] (s2.north) |- (ml.west);
		\node[label, right=0.1cm of ml] {sanity check\\$(C_{11},C_{12})$ predictor};
	\end{tikzpicture}
	\caption{Five-stage pipeline. Stages~1 through 4 seed the potential and
	Stage~5b refits it. The dashed branch above Stage~3 is the composition
	surrogate of Section~\ref{sec:comp_nn}, which reads the Stage~2 elastic
	database as a consistency check and never feeds back into the
	potential-fitting flow.}
	\label{fig:pipeline}
\end{figure}

The rest of Section~\ref{sec:progress} reports each stage in turn: what it
does, what it found, the problems that came up, and how they were handled.
Stages~1 through 4 are covered in
Sections~\ref{sec:stage1}--\ref{sec:stage4}, the two Stage~5 branches in
Sections~\ref{sec:comp_nn} and~\ref{sec:nnip}, and the external validation in
Section~\ref{sec:validation}.

\newpage
